% TOP_PROBLEM: {"problemId": "problem.combinatorial-invariance-of-kazhdan-lusztig-polynomials", "canonicalTitle": "Combinatorial Invariance of Kazhdan-Lusztig Polynomials", "releaseRank": 423, "primaryDomain": "algebra_representation_category", "primaryDomainLabel": "Algebra, representation & category theory", "displayStatus": "open_with_solved_subcases", "releaseStatus": "open_with_solved_subcases"}
% !TeX program = lualatex
% ATTEMPT_STATUS: unresolved
% Research continuation by GPT_6_Astra_Ultra, 27 September 2026.
\documentclass[11pt]{article}
\usepackage[margin=25mm]{geometry}
\usepackage{fontspec}
\setmainfont{Latin Modern Roman}
\usepackage{amsmath,amssymb,mathtools}
\usepackage{unicode-math}
\setmathfont{Latin Modern Math}
\usepackage{xurl,hyperref}
\hypersetup{colorlinks=true,urlcolor=blue}
\setcounter{secnumdepth}{0}
\setlength{\parindent}{0pt}
\setlength{\parskip}{5pt}
\setlength{\emergencystretch}{3em}
\begin{document}
\section*{\texorpdfstring{Combinatorial Invariance of Kazhdan-Lusztig Polynomials}{Combinatorial Invariance of Kazhdan-Lusztig Polynomials}}\label{423}
\subsection{Definitions and mathematical statement}
A Coxeter system $(W,S)$ has generators $s^2=1$ and relations $(st)^{m_{st}}=1$; let $\ell$ be word length and $T$ the set of conjugates of simple reflections. Bruhat order is the subword order on reduced expressions. On an interval $[u,v]$, use the directed Bruhat graph with edges $x\to y$ whenever $yx^{-1}\in T$ and $\ell(x)<\ell(y)$, not just length-one covers. The Kazhdan--Lusztig polynomials $P_{x,y}(q)$ are the standard triangular coefficients of the bar-invariant basis of the Coxeter Hecke algebra, normalized by $P_{x,x}=1$ and $\deg P_{x,y}\le(\ell(y)-\ell(x)-1)/2$ for $x<y$. The conjecture says a directed-graph isomorphism $[u,v]\cong[w,z]$, even across different Coxeter systems, forces $P_{u,v}=P_{w,z}$. Reflection labels on edges are not supplied to the graph isomorphism.
\par\smallskip\textit{Formulation and concept references:} \href{https://link.springer.com/article/10.1007/s00029-024-00972-0}{[S1]}.

\subsection{Short English statement}
Does the unlabelled directed Bruhat graph of an interval determine its Kazhdan--Lusztig polynomial, independently of the surrounding Coxeter group?
\subsection{Research attempt}
\textit{GPT-6 Astra Ultra; 27 September 2026. Status: UNRESOLVED.}

I extract all the data needed for intervals of length at most three
from the unlabelled directed reflection graph. This gives an explicit
formula in length three and a finite test of its normalization. The
calculation also identifies the obstacle to continuing the same argument:
the known general path formula uses a reflection ordering, information
which the proposed graph isomorphism does not supply.

\paragraph{Rank is intrinsic to the directed graph.}
Every directed edge strictly increases Coxeter length. A directed path
from $u$ to $x$ therefore has at most $\ell(x)-\ell(u)$ edges.
On the other hand a saturated Bruhat chain from $u$ to $x$ has exactly
that many edges, all of which are reflection edges. It follows that
\[
 r(x)=\ell(x)-\ell(u)
 =\max\{\text{number of edges in a directed path from }u\text{ to }x\}.
                                                        \tag{1}
\]
The unique source and sink identify the interval endpoints. Thus graph
isomorphisms preserve rank, length difference, atoms, and coatoms.
Cover edges are precisely the edges whose ranks differ by one.
The converse reconstruction is not immediate: the reflection graph
may also contain longer edges, which must be retained.

For rank zero, $P_{u,u}=1$. For ranks one and two the degree bound and
the standard constant-term identity $P_{u,v}(0)=1$ give $P_{u,v}=1$.
These standard elementary properties of the Hecke normalization are
used throughout, together with the bar relation
\[
 q^dP_{u,v}(q^{-1})=
 \sum_{u\le x\le v} R_{u,x}(q)P_{x,v}(q),
 \qquad d=\ell(v)-\ell(u).                           \tag{2}
\]
Here $R_{x,x}=1$, rank-one $R$-polynomials equal $q-1$, and
rank-two ones equal $(q-1)^2$. The general recursion determining these
polynomials is recalled in [E1].

\paragraph{An explicit graph formula in rank three.}
Assume $d=3$. Let $a$ be the number of atoms of the interval and let
$\varepsilon$ equal one if there is a directed edge $u\to v$, and
zero otherwise. I claim
\[
 P_{u,v}(q)=1+(a-3+\varepsilon)q.                     \tag{3}
\]
To prove it, use Dyer's increasing-path formula [E1] for normalized
$R$-polynomials. In its convention
\[
 R_{u,v}(q)=q^{d/2}\widetilde R_{u,v}
                         (q^{1/2}-q^{-1/2}).          \tag{4}
\]
The coefficient of $t^j$ in $\widetilde R$ counts paths of $j$
reflection edges whose labels increase in a reflection order. Each
reflection changes the parity of length. A path with total length
difference three therefore has either one or three edges. The one-edge
path exists exactly when $\varepsilon=1$, and its label is automatically
increasing. The leading coefficient of the monic degree-three
$R$-polynomial gives coefficient one on $t^3$. Hence
\[
 \widetilde R_{u,v}(t)=t^3+\varepsilon t,
\]
and (4) yields
\[
 R_{u,v}(q)=(q-1)^3+\varepsilon q(q-1)
 =q^3+(\varepsilon-3)q^2+(3-\varepsilon)q-1.          \tag{5}
\]
The path formula and monicity are established inputs; the specialization
to this low rank is the calculation here.

Bruhat intervals are Eulerian, so the number of coatoms in rank three
also equals $a$: the alternating rank sum is $1-a+b-1=0$.
Write $P_{u,v}=1+cq$. Every proper subinterval appearing in (2) has
length at most two and thus has $P=1$. Separating the four ranks in
(2) gives
\[
 q^3+cq^2=(1+cq)+a(q-1)+a(q-1)^2+R_{u,v}(q).
\]
Substitution of (5), followed by comparison of the coefficient of $q$,
gives $c=a-3+\varepsilon$. This proves (3) for any Coxeter system,
not just a symmetric group. Each quantity on the right side is visible
in the unlabelled directed graph. The selected invariance assertion
therefore holds for all intervals of length at most three.

\paragraph{Checks distinguishing reflection edges from covers.}
In type $A_2$, the interval from the identity to the longest element
has rank three and two atoms. The longest element is itself a
reflection, so $\varepsilon=1$ and (3) gives $P=1$. Dropping the
long edge would instead give the incorrect coefficient $-1$.
A rank-three Boolean interval has three atoms and no endpoint
reflection edge, again giving $P=1$. These tests check both the sign
and the need for the graph convention in the question.

An exact Python calculation enumerates all 24 permutations of $S_4$.
It tests Bruhat order by sorted-prefix inequalities and computes
$R$-polynomials by the usual simple-descent recurrence. It then
reconstructs $P$ from (2) and the degree bound, checking the entire
polynomial identity rather than just its low coefficients. There are
213 comparable intervals in this enumeration, including 42 of rank
three; all satisfy (3), and two rank-three intervals have a nonconstant
polynomial. This is a reproducible finite check of the formula and
conventions. It is not a proof in other ranks or other Coxeter groups;
the preceding symbolic argument supplies the stated general low-rank
result.

\paragraph{Why the general path argument stops.}
For larger rank, increasing paths of intermediate lengths contribute
to $\widetilde R$. Their number is not the total number of directed
paths of that length. Whether their successive reflections increase
depends on an ordering of reflection labels. The unlabelled graph
retains endpoints and adjacency, but an arbitrary ordering of its
edges need not arise from a reflection order and has no justification
in Dyer's formula. Using the formula with invented labels would not
prove invariance under the isomorphisms in the question.

Likewise the familiar descent recurrence compares $u,v$ with $us,vs$
for a simple reflection $s$. The ambient system supplies that operation;
the interval graph alone does not come with a distinguished $s$ or a
map implementing those multiplications. A valid recursive proof must
either reconstruct sufficient descent data intrinsically or show that
all allowed choices produce the same result without those data.

\paragraph{Source scope and outcome.}
[S1] discusses precisely the directed-graph formulation, which is
retained here. The primary path formula [E1] provides a tool with
additional labels, not the missing unlabelled invariance theorem.
The first unresolved step in this route is a graph-intrinsic way to
recover the intermediate increasing-path contributions, or an
alternative intrinsic recurrence. The rank-three identity (3) and its
finite checks are the strongest conclusions of the attempt. The full
combinatorial invariance conjecture remains \textbf{UNRESOLVED}.

\subsection{Sources}
\begin{itemize}
\item[S1]Gurevich and Wang, Parabolic recursions for Kazhdan-Lusztig polynomials and the hypercube decomposition, Selecta Mathematica (2024), Conjecture 5.1.
\url{https://link.springer.com/article/10.1007/s00029-024-00972-0}.
\textit{Formulation source. Consulted 24 September 2026: Full directed reflection Bruhat graphs and their combinatorial invariance conjecture.}
\item[E1]Matthew J. Dyer, \emph{Hecke algebras and shellings of Bruhat
intervals}, Compositio Mathematica 89 (1993), 91--115.
\url{https://www.numdam.org/item/CM_1993__89_1_91_0.pdf}.
\textit{Primary source for the reflection-order increasing-path formula
and the Hecke-algebra conventions used in the rank-three calculation.
The labels in that formula are additional data, not supplied by the
conjecture's graph isomorphism. Consulted 27 September 2026.}

\end{itemize}
\end{document}
